% SPDX-License-Identifier: MIT-0
\documentclass[11pt]{article}
\usepackage[T1]{fontenc}
\usepackage[utf8]{inputenc}
\usepackage{lmodern}
\usepackage[a4paper,margin=25mm]{geometry}
\usepackage{amsmath,amssymb,amsthm,mathtools}
\usepackage{booktabs,longtable,array,multicol}
\usepackage{microtype,enumitem,fancyhdr}
\usepackage[dvipsnames]{xcolor}
\usepackage{listings}
\usepackage{hyperref}
\definecolor{Ink}{HTML}{203545}
\definecolor{Accent}{HTML}{536348}
\hypersetup{colorlinks=true,linkcolor=Ink,citecolor=Accent,urlcolor=Ink,
 pdftitle={Coefficientwise magic positivity through dimension ten},
 pdfauthor={Research draft prepared with ChatGPT for Vladimir Reshetnikov},
 pdfsubject={A computer-assisted partial solution of a generalized parking-function polytope conjecture}}
\newtheorem{theorem}{Theorem}[section]
\newtheorem{proposition}[theorem]{Proposition}
\newtheorem{lemma}[theorem]{Lemma}
\newtheorem{corollary}[theorem]{Corollary}
\theoremstyle{definition}
\newtheorem{definition}[theorem]{Definition}
\newtheorem{conjecture}[theorem]{Conjecture}
\newtheorem{example}[theorem]{Example}
\newtheorem{remark}[theorem]{Remark}
\newcommand{\N}{\mathbb N}
\newcommand{\Z}{\mathbb Z}
\newcommand{\Q}{\mathbb Q}
\newcommand{\R}{\mathbb R}
\newcommand{\bb}{\mathbf b}
\newcommand{\xx}{\mathbf x}
\newcommand{\qq}{\mathbf q}
\newcommand{\one}{\mathbf 1}
\newcommand{\zero}{\mathbf 0}
\newcommand{\X}{\mathfrak X}
\newcommand{\M}{\mathcal M}
\newcommand{\supp}{\operatorname{supp}}
\newcommand{\interior}{\operatorname{int}}
\newcommand{\conv}{\operatorname{conv}}
\newcommand{\degree}{\operatorname{deg}}
\DeclareMathOperator{\proj}{proj}
\numberwithin{equation}{section}
\pagestyle{fancy}
\fancyhf{}
\fancyhead[L]{\small\color{Ink}Magic positivity through dimension ten}
\fancyhead[R]{\small Research draft}
\fancyfoot[C]{\thepage}
\setlength{\headheight}{14pt}
\setlength{\parskip}{0.35em}
\emergencystretch=2em
\lstset{basicstyle=\small\ttfamily,breaklines=true,columns=fullflexible,
 frame=single,rulecolor=\color{black!20},showstringspaces=false,
 xleftmargin=0.5em,xrightmargin=0.5em}
\title{\LARGE\bfseries\color{Ink}Coefficientwise magic positivity\\through dimension ten\\[0.45em]
\large An exact-arithmetic partial solution for\\generalized parking-function polytopes}
\author{Research draft prepared for Vladimir Reshetnikov}
\date{20 September 2026}

\begin{document}
\maketitle
\begin{abstract}
Hill, Luo, Trinh, and Vindas-Mel\'endez conjecture that every generalized
parking-function polytope $\X_n(\bb)$ with $n\ge3$ is magic positive: its Ehrhart
polynomial has nonnegative coefficients in the basis
$t^j(t+1)^{n-j}$. We obtain an exact-arithmetic, computer-assisted proof for
\emph{all positive integer parameter vectors} in every dimension $3\le n\le10$.
The certified statement is stronger: after $b_i=1+x_i$, every coefficient of
every polynomial $n!\mu_{n,j}(\one+\xx)$ is nonnegative. For $4\le n\le10$,
all monomials of total degree at most $j$ occur with strictly positive
coefficients. Dimension three has a short explicit proof printed in full.
Independently of the dimension bound, we derive a harmonic-number formula for
$\mu_{n,1}$ and prove its sharp minimum $n(3n-7)/4$. We also identify the zero
cases of the final magic coefficient and give a cube-addition identity reducing
the general problem to $b_1=1$. The accompanying package contains exact
polynomial certificates, a standard-library Python generator and verifier,
and independent lattice-point checks. The unrestricted conjecture is not
settled here; the distinction between the finite-dimensional theorem and the
remaining all-dimensional problem is maintained throughout.
\end{abstract}

\noindent\textbf{Keywords.} Ehrhart polynomial; magic positivity; generalized
parking function; polymatroid; exact polynomial certificate.

\setcounter{tocdepth}{1}
\begingroup
\setlength{\parskip}{0pt}
\small
\tableofcontents
\endgroup
\clearpage

\section{The selected problem and the result obtained}
\label{sec:problem}

\subsection{Selection, exclusion, and current source status}
The subject was selected before choosing the problem. A catalogue of 120
mathematical areas was fixed, and a pseudorandom draw seeded by 128 bits from the
operating system selected entry 33, \emph{discrete geometry}. There was no
redraw. The seed, ordered catalogue, and replay command are recorded in
Appendix~\ref{app:selection} and in \texttt{data/selection.json}.

The supplied \texttt{manifest(1).tex} was read in full and used as an exclusion
list. Its 71 entries do not contain the problem considered here. In particular,
its lucky-parking-spots problem concerns a statistic on parking functions;
its preorder-polytope and order/chain-polytope problems concern different
polytopes and different invariants. The present question concerns the
Ehrhart polynomial of the convex hull of generalized parking functions.
The input-file hash and the exclusion comparison are recorded in
\texttt{notes/exclusion\_audit.md}.

The precise target is Conjecture~8.1 of Hill, Luo, Trinh, and
Vindas-Mel\'endez \cite{HLTV}. Their preprint, arXiv:2607.15503v1, was
submitted on 16 July 2026. That version explicitly leaves arbitrary parameter vectors open,
while proving the two-parameter family $(a,b,\ldots,b)$. The version history
and targeted literature searches were checked on 20 September 2026; no later
resolution of the arbitrary-vector conjecture was located. The investigation below uses that published formulation.

\subsection{Definitions and the conjecture}
For $\bb=(b_1,\ldots,b_n)\in\Z_{>0}^n$, put
\[
 S_i=b_1+\cdots+b_i,\qquad S_0=0.
\]
A \emph{$\bb$-parking function} is a positive integer vector whose increasing
rearrangement $\beta'_1\le\cdots\le\beta'_n$ satisfies $\beta'_i\le S_i$.
Let $\X_n(\bb)$ be the convex hull of these vectors, and define
\[
 E_{n,\bb}(t)=\#\bigl(t\X_n(\bb)\cap\Z^n\bigr)
 \qquad(t\in\Z_{\ge0}).
\]
The vertices are integral, and this counting function is a polynomial. We
will also obtain its polynomiality directly from an exact recursion.

For $n\ge2$, these polytopes have dimension $n$. Expand uniquely as
\begin{equation}
 E_{n,\bb}(t)=\sum_{j=0}^n\mu_{n,j}(\bb)t^j(t+1)^{n-j}.
 \label{eq:magic-basis}
\end{equation}
The word \emph{magic} names the basis, not an additional geometric operation.
The polytope is \emph{magic positive} when all $\mu_{n,j}(\bb)$ are
nonnegative. Equivalently, the polynomial
\begin{equation}
 M_{n,\bb}(y)=(1-y)^n E_{n,\bb}\!\left(\frac{y}{1-y}\right)
       =\sum_{j=0}^n\mu_{n,j}(\bb)y^j
 \label{eq:transform}
\end{equation}
has nonnegative ordinary coefficients.

\begin{conjecture}[Hill--Luo--Trinh--Vindas-Mel\'endez]
\label{conj:target}
For every $n\ge3$ and every $\bb\in\Z_{>0}^n$, the polytope $\X_n(\bb)$ is
magic positive.
\end{conjecture}

Magic positivity is stronger than ordinary Ehrhart positivity. It also
implies real-rootedness of the $h^*$-polynomial through Br\"and\'en's
transformation theorem \cite[Theorem~4.2]{Branden}. This explains why changing
the basis is useful rather than merely cosmetic. Section~\ref{sec:consequences}
makes the implication precise.

\subsection{Main results and their proof status}
\begin{theorem}[Coefficientwise positivity through dimension ten]
\label{thm:main}
For $3\le n\le10$ and $0\le j\le n$,
\begin{equation}
 n!\mu_{n,j}(1+x_1,\ldots,1+x_n)\in
       \Z_{\ge0}[x_1,\ldots,x_n].
 \label{eq:strong}
\end{equation}
For $4\le n\le10$, every monomial of total degree at most $j$ has a
strictly positive coefficient. For $n=3$, the same assertion holds except
that the constant coefficient of $6\mu_{3,3}(\one+\xx)$ is zero.
Consequently Conjecture~\ref{conj:target} holds for all positive integer
parameter vectors in these eight dimensions.
\end{theorem}

The dimension-three proof is entirely explicit in
Section~\ref{sec:cubic}. The proof for dimensions four through ten is
computer-assisted: exact polynomial identities obtained from a proved
recursion are checked, and their complete coefficient lists are supplied.
It is not a numerical test over a finite box of parameter values.

\begin{theorem}[The first magic coefficient in every dimension]
\label{thm:first-intro}
Let $H_m=\sum_{k=1}^m1/k$. For $n\ge3$ and $x_i=b_i-1\ge0$,
\begin{equation}
 \mu_{n,1}(\one+\xx)
 =\frac{n(3n-7)}4+n x_1+
  \sum_{r=1}^{n-1}r(1+H_n-H_r)x_{n-r+1}.
 \label{eq:first-intro}
\end{equation}
In particular,
\[
 \mu_{n,1}(\bb)\ge \frac{n(3n-7)}4>0,
\]
with equality precisely at $\bb=\one$.
\end{theorem}

The second theorem has a proof in arbitrary dimension and does not rely on
the certificates. The identities $\mu_{n,0}=1$ and
$\mu_{n,n}=\#(\interior\X_n(\bb)\cap\Z^n)$ further control the endpoints.
Therefore any counterexample not ruled out by this work must have
$n\ge11$ and a negative coefficient at an index $2\le j\le n-1$.
This does \emph{not} assert that such a counterexample exists.

The conjecture in all dimensions, and the coefficientwise strengthening of
it, remain unproved in this draft.

\section{Geometry and an exact counting recursion}
\label{sec:geometry}

The inequality description used here is due to Bayer and collaborators
\cite[Theorem~2.3]{Bayer}. The lattice-layer recursion, polynomiality, and
dilation identity are established in \cite[Theorem~3.5, Theorem~4.1,
Proposition~5.4, and Section~5]{HLTV}. They are inputs rather than novelty
claims. We provide direct derivations because the correctness of the
certificates depends on their boundary conventions.

\subsection{A symmetric rank description}
Translate by the integer vector $\one$, and write
\[
 P_n(\bb)=\X_n(\bb)-\one,\qquad u_i=S_i-1,
 \qquad f(k)=\sum_{i=n-k+1}^n u_i,\quad f(0)=0.
\]
Then
\begin{equation}
 P_n(\bb)=\left\{z\in\R_{\ge0}^n:
       \sum_{i\in I}z_i\le f(|I|)\text{ for every }I\subseteq[n]\right\}.
 \label{eq:rank-polytope}
\end{equation}
The sequence $u_1\le\cdots\le u_n$ is nonnegative. The same definition makes
sense for the auxiliary weak parameter range
\begin{equation}
 b_1\ge1,\qquad b_2,\ldots,b_n\ge0.
 \label{eq:weak}
\end{equation}
We will need this range to evaluate the counting polynomial on coordinate
axes. It is a genuine geometric extension, not an assumption that arbitrary
polynomial specializations still count lattice points.

\begin{lemma}[Greedy description]
\label{lem:greedy}
The right side of \eqref{eq:rank-polytope} is the convex hull of all coordinate
permutations of
\[
 (u_n,u_{n-1},\ldots,u_{n-m+1},0,\ldots,0),\qquad 0\le m\le n.
\]
This holds throughout \eqref{eq:weak}. In particular the polytope is integral
when the parameters are integers.
\end{lemma}
\begin{proof}
Each displayed vector satisfies all subset inequalities, since the sum of
any $k$ coordinates is at most the sum of the $k$ largest $u_i$.
For the reverse inclusion, maximize a linear functional. By symmetry, its
positive coefficients may be ordered $c_1\ge\cdots\ge c_m>0$; coordinates
with nonpositive coefficients can be set to zero, since the feasible set is
downward closed. With $c_{m+1}=0$, every feasible $z$ satisfies
\[
 \sum_{i=1}^m c_i z_i
 =\sum_{k=1}^m(c_k-c_{k+1})\sum_{i=1}^kz_i
 \le\sum_{k=1}^m(c_k-c_{k+1})f(k).
\]
Equality is attained by $z_i=u_{n-i+1}$ for $i\le m$ and zero otherwise.
Thus the proposed convex hull has the same support function. After adding
$\one$, the displayed vectors are parking functions, so this also identifies
the original convex hull when the parameters are positive.
\end{proof}

\begin{remark}
Counting lattice points of the convex hull is not the same as counting
parking functions. For example, $(2,2,2)\in\X_3(1,1,1)$, since its largest
one-, two-, and three-coordinate sums are $2,4,6$, bounded by $3,5,6$.
It is not a parking function: its smallest coordinate exceeds $S_1=1$.
The algorithms below count points of the polytope, not just its defining
integer set.
\end{remark}

\subsection{Derivation of the layers}
In the unshifted coordinates put
\[
 \Sigma_k=\sum_{i=n-k+1}^n S_i.
\]
Fix the last coordinate at an integer height $h$, where $1\le h\le S_n$.
For a $k$-element subset $I$ of the other coordinates, the two relevant
bounds are
\[
 \sum_{i\in I}x_i\le\Sigma_k,
 \qquad \sum_{i\in I}x_i\le\Sigma_{k+1}-h.
\]
Their minimum is the exact residual bound. Indeed all inequalities either
contain the fixed coordinate or do not. The switch between the two bounds
is controlled by
\[
 \Sigma_{k+1}-\Sigma_k=S_{n-k}.
\]

Write $h=S_{\ell-1}+r$ with $1\le r\le b_\ell$. The resulting layer is
$\X_{n-1}(\bb^{(\ell,r)})$, where
\begin{align}
 \bb^{(1,r)}&=(b_1+b_2,b_3,\ldots,b_n),\label{eq:first-layer}\\
 \bb^{(\ell,r)}&=(b_1,\ldots,b_{\ell-2},
 b_{\ell-1}+b_\ell-r,b_{\ell+1}+r,b_{\ell+2},\ldots,b_n),
 &&2\le\ell\le n-1,\label{eq:middle-layer}\\
 \bb^{(n,r)}&=(b_1,\ldots,b_{n-2},b_{n-1}+b_n-r).
 \label{eq:last-layer}
\end{align}
To check the indexing, the new partial sums for $\ell>1$ are
$S_i$ for $i\le\ell-2$, then $S_\ell-r$, and then $S_{i+1}$ for
$i\ge\ell$; for $\ell=1$ they are simply $S_{i+1}$.
Their largest-coordinate sums give precisely
$\min(\Sigma_k,\Sigma_{k+1}-h)$.

This derivation also covers zero values among $b_2,\ldots,b_n$: a block
of length zero contributes no layers. Every parameter vector that does
occur still satisfies \eqref{eq:weak}.

\begin{proposition}[Exact lattice-count recurrence]
\label{prop:recurrence}
Let $L_n(\bb)=\#(\X_n(\bb)\cap\Z^n)$ in the weak integer parameter range.
Then $L_1(b_1)=b_1$, and for $n\ge2$,
\begin{align}
 L_n(\bb)={}&b_1L_{n-1}(b_1+b_2,b_3,\ldots,b_n)\notag\\
 &+\sum_{\ell=2}^{n-1}\sum_{r=1}^{b_\ell}
       L_{n-1}(\bb^{(\ell,r)})
   +\sum_{r=1}^{b_n}L_{n-1}(\bb^{(n,r)}).
 \label{eq:recurrence}
\end{align}
The middle sum is empty for $n=2$.
\end{proposition}
\begin{proof}
The layers partition the lattice points. In the first block all $b_1$ layers
have the same parameter vector \eqref{eq:first-layer}; the other blocks use
\eqref{eq:middle-layer} and \eqref{eq:last-layer} individually.
\end{proof}

\subsection{Polynomiality, degree, and denominators}
\begin{proposition}
\label{prop:polynomial}
There is a unique polynomial $L_n\in\Q[b_1,\ldots,b_n]$ of total degree at most
$n$ that gives these counts. Moreover,
\begin{equation}
 n!L_n\in\Z[b_1,\ldots,b_n].
 \label{eq:integral-scale}
\end{equation}
\end{proposition}
\begin{proof}
Induct on $n$. Affine substitution preserves the degree bound $n-1$.
Expanding a summand in powers of $r$ and summing $r^k$ from $1$ to $b_\ell$
replaces it by a polynomial of degree $k+1$; hence the total degree rises by
at most one. The first block also raises degree by at most one. This
constructs $L_n$. Uniqueness follows because a polynomial vanishing on all
positive integer vectors vanishes identically, by induction on the number
of variables.

For the denominator statement, let $R(\xx)=L_n(\one+\xx)$. Its values on
$\Z_{\ge0}^n$ are integers. The multivariate Newton expansion is
\[
 R(\xx)=\sum_{|\alpha|\le n}
 (\Delta_1^{\alpha_1}\cdots\Delta_n^{\alpha_n}R)(\zero)
 \prod_{i=1}^n\binom{x_i}{\alpha_i}.
\]
All finite differences are integers. Each denominator
$\prod_i\alpha_i!$ divides $n!$, because $|\alpha|\le n$.
Multiplication by $n!$ therefore gives integral ordinary coefficients;
translation back to the $b_i$ preserves integrality.
\end{proof}

The recurrence and the polynomial agree on the weak parameter domain,
not merely on its strictly positive part. This follows from the layer proof,
where zero-length blocks are omitted. This point is useful in the next section
and in the all-dimensional coefficient calculation.

\section{The magic transform and finite symbolic certificates}
\label{sec:certificates}

\subsection{Dilation and homogeneous components}
By the rank description, multiplication of $P_n(\bb)$ by a nonnegative
integer $t$ replaces $u_i$ by $t u_i$. Consequently,
\begin{equation}
 E_{n,\bb}(t)=L_n(1+t(b_1-1),t b_2,\ldots,t b_n).
 \label{eq:dilation}
\end{equation}
The translation between $t\X_n(\bb)$ and $tP_n(\bb)$ is integral.
At $t=0$, the right side is $L_n(1,0,\ldots,0)=1$, as required.
Thus \eqref{eq:dilation} supplies an Ehrhart polynomial without an
interpolation step.

Define
\[
 F_n(\qq)=L_n(1+q_1,q_2,\ldots,q_n)
     =\sum_{d=0}^n F_{n,d}(\qq),
\]
where $F_{n,d}$ is homogeneous of degree $d$. If
$E_{n,\bb}(t)=\sum_{d=0}^n e_{n,d}(\bb)t^d$, then
\begin{equation}
 e_{n,d}(\bb)=F_{n,d}(b_1-1,b_2,\ldots,b_n),\qquad
 \mu_{n,j}=\sum_{d=0}^j(-1)^{j-d}\binom{n-d}{j-d}e_{n,d}.
 \label{eq:triangular-transform}
\end{equation}
In particular $\mu_{n,j}(\one+\xx)$ has degree at most $j$.
The integer scaling in \eqref{eq:integral-scale} survives all translations
and the integral triangular transform.

There are two different translations here. The homogeneous grading uses
$(q_1,q_2,\ldots)=(b_1-1,b_2,\ldots)$, whereas the positivity certificate
uses $x_i=b_i-1$ for \emph{every} coordinate. Confusing these shifts would
produce an incorrect certificate; the implementation performs them in
separate stages.

\subsection{A small exact summation kernel}
For nonnegative integers $a,b$, introduce
\begin{equation}
 K_{a,b}(X,Y,Z)=\sum_{r=1}^{Y}(X+Y-r)^a(Z+r)^b,
 \label{eq:kernel}
\end{equation}
where the sum is interpreted as its unique polynomial continuation in $Y$.
This kernel carries every middle block of the recursion. The final block
uses $K_{a,0}(X,Y,0)$.

For a complete construction, let $A_k(Y)=\sum_{r=1}^Yr^k$. Telescoping
$(r+1)^{k+1}-r^{k+1}$ gives
\begin{equation}
 A_k(Y)=\frac{(Y+1)^{k+1}-1-
               \sum_{h=0}^{k-1}\binom{k+1}{h}A_h(Y)}{k+1}.
 \label{eq:power-sum}
\end{equation}
After binomial expansion,
\begin{equation}
 K_{a,b}(X,Y,Z)=
 \sum_{i+j+h=a}\binom{a}{i,j,h}X^iY^j(-1)^h
 \sum_{z=0}^{b}\binom bz Z^z A_{h+b-z}(Y).
 \label{eq:kernel-expansion}
\end{equation}
All arithmetic is rational, and $n!K_{a,b}$ has integral coefficients when
$a+b+1\le n$. This also follows from the integer-valued-polynomial argument
of Proposition~\ref{prop:polynomial}.

An independent identity is especially useful for auditing the kernel:
\begin{align}
 K_{a,b}(X,0,Z)&=0,\label{eq:kernel-boundary}\\
 K_{a,b}(X,Y,Z)-K_{a,b}(X+1,Y-1,Z)
     &=X^a(Z+Y)^b.\label{eq:kernel-difference}
\end{align}
The second term has exactly the first $Y-1$ summands of the first; the
remaining summand is the displayed right side. Conversely these identities
uniquely determine the kernel on nonnegative integer $Y$, and therefore
as a polynomial. The verifier checks them by independent polynomial
translation, not by reapplying the power-sum recurrence.

\subsection{What the certificate proves}
A sparse polynomial is represented by a map
\[
 (\alpha_1,\ldots,\alpha_n)\longmapsto c_\alpha\in\Z.
\]
The lattice files store $n!L_n(\bb)$, and the magic files store
$n!\mu_{n,j}(\one+\xx)$. Missing monomials have coefficient zero.
The three invariants of the generator are:
\begin{enumerate}[label=(\roman*),itemsep=0.15em]
\item the current lattice map is exactly $n!L_n$, by the layer recursion;
\item homogeneous extraction and \eqref{eq:triangular-transform} produce
exactly the scaled magic maps;
\item a nonnegative map evaluates nonnegatively whenever $x_i\ge0$.
\end{enumerate}
No root approximation, positivity tolerance, guessed recurrence, or
parameter-bound assumption enters any invariant.

\begin{proof}[Computer-assisted proof of Theorem~\ref{thm:main}]
Start from $L_1(b_1)=b_1$. Equations \eqref{eq:recurrence},
\eqref{eq:power-sum}, and \eqref{eq:kernel-expansion} construct $L_n$
inductively. Equations \eqref{eq:dilation} and
\eqref{eq:triangular-transform} construct the magic maps. The supplied
verifier regenerates and compares every stored coefficient for $1\le n\le10$.
For $3\le n\le10$, every retained coefficient is a positive integer.
The number of retained monomials in the $j$th map is
\begin{equation}
 \binom{n+j}{j}-\begin{cases}1,&(n,j)=(3,3),\\0,&\text{otherwise}.
 \end{cases}
 \label{eq:support-size}
\end{equation}
The degree bound is $j$, so for $n\ge4$ this is the entire possible support.
For $n=3,j=3$, the single missing monomial is the constant, as the explicit
formula below confirms. Evaluating at $x_i=b_i-1\ge0$ proves the stated
infinite-parameter assertions. The exact output is summarized in
Table~\ref{tab:certificate-counts}; the full coefficient lists, rather than
only their counts, are included in the archive.
\end{proof}

\begin{table}[htbp]
\centering
\begin{tabular}{@{}rrrr@{}}
\toprule
$n$ & Terms in $n!L_n$ & Terms in all magic maps & $\mu_{n,n}(\one)$\\
\midrule
3 & 16 & 34 & 0 \\
4 & 60 & 126 & 5 \\
5 & 227 & 462 & 96 \\
6 & 851 & 1,716 & 1,758 \\
7 & 3,223 & 6,435 & 34,820 \\
8 & 12,200 & 24,310 & 762,675 \\
9 & 46,456 & 92,378 & 18,489,870 \\
10 & 177,525 & 352,716 & 493,794,476 \\
\bottomrule
\end{tabular}
\caption{Exact certificate sizes and the final magic coefficient at the
classical parameter vector. All 478,177 retained coefficients in the magic
maps for $3\le n\le10$ are positive. The last column is also an interior
lattice-point count.}
\label{tab:certificate-counts}
\end{table}

The number of possible monomials in a polynomial of degree at most $j$ in
$n$ variables is $\binom{n+j}{j}$. Summing over $j$ gives
$\binom{2n+1}{n}$, which explains the total support size for $n\ge4$.
These support counts are observations proved for the listed dimensions,
not an all-dimensional extrapolation.

\section{A complete explicit proof in dimension three}
\label{sec:cubic}

The first two recursion steps give
\begin{align*}
 L_2(b_1,b_2)&=b_1^2+2b_1b_2+\tfrac12b_2^2-\tfrac12b_2,\\
 L_3(b_1,b_2,b_3)&=b_1^3+3b_1^2b_2+3b_1^2b_3
 +3b_1b_2^2+6b_1b_2b_3+\tfrac32b_1b_3^2-\tfrac32b_1b_3\\
 &\quad+\tfrac56b_2^3+2b_2^2b_3-\tfrac12b_2^2
 +b_2b_3^2-2b_2b_3-\tfrac13b_2
 +\tfrac16b_3^3-\tfrac12b_3^2+\tfrac13b_3.
\end{align*}
The $L_2$ formula follows by summing the last block of
\eqref{eq:recurrence}; substitution of it into the same recurrence gives
$L_3$ using sums of first and second powers. Thus the following identities
can be checked using only elementary polynomial arithmetic.

Put $x=b_1-1$, $y=b_2-1$, and $z=b_3-1$. Applying
\eqref{eq:triangular-transform} gives
\begin{align}
 6\mu_{3,0}&=6,\notag\\
 6\mu_{3,1}&=18x+16y+11z+9,\label{eq:cubic-mu1}\\
 6\mu_{3,2}&=18x^2+x(36y+27z+27)
       +15y^2+24yz+22y+6z^2+14z+9,\label{eq:cubic-mu2}\\
 6\mu_{3,3}&=6x^3+18x^2(y+z+1)\notag\\
 &\quad+x(18y^2+36yz+36y+9z^2+27z+18)\notag\\
 &\quad+5y^3+12y^2z+12y^2+6yz^2+12yz+7y
       +z^3+3z^2+2z.\label{eq:cubic-mu3}
\end{align}
All coefficients are nonnegative, and $x,y,z\ge0$ for the desired
parameter range. This proves Conjecture~\ref{conj:target} in dimension
three without any unprinted computation. Moreover, $\mu_{3,1}$ and
$\mu_{3,2}$ are always positive, while $\mu_{3,3}=0$ holds exactly when
$x=y=z=0$.

\begin{example}[An asymmetric parameter vector]
For $\bb=(2,3,1)$, the formulas give
\begin{align*}
 M_{3,\bb}(y)&=1+\frac{59}{6}y+\frac{115}{3}y^2+54y^3,\\
 E_{3,\bb}(t)&=1+\frac{77}{6}t+61t^2+\frac{619}{6}t^3,\\
 h^*_{3,\bb}(z)&=1+174z+390z^2+54z^3.
\end{align*}
In particular there are 178 lattice points at dilation one. The first line
agrees with the independently printed example in \cite[Section~8]{HLTV};
that example was used as a check, not to fit the polynomial.
\end{example}

\subsection{Why the lower-dimensional exception is genuine}
For $u=b_1-1$ and $v=b_2-1$, the degree-two magic coefficients are
\[
 2\mu_{2,0}=2,\qquad 2\mu_{2,1}=4u+3v-1,\qquad
 2\mu_{2,2}=2u^2+4uv+v(v-1).
\]
For nonnegative integers $u,v$, the final expression is nonnegative; the
middle one is negative precisely when $u=v=0$. This rederives the known
exception $\X_2(1,1)$ and shows why a proof that assumes coefficientwise
positivity beginning in dimension two would fail. The nonnegative-value
statement and the nonnegative-ordinary-coefficient statement are distinct:
$v(v-1)$ has a negative linear coefficient but is nonnegative at every
nonnegative integer.

In dimension one, $\X_1(b_1)=[1,b_1]$ has Ehrhart polynomial
$1+(b_1-1)t$. For $b_1\ge2$ its degree-one magic transform is
$1+(b_1-2)y$. The case $b_1=1$ is a point of intrinsic dimension zero;
its intrinsic magic polynomial is $1$. The auxiliary certificate files
use ambient degree one in that case, yielding $1-y$; this is not a
counterexample to the intrinsic-dimensional definition.

\section{A sharp all-dimensional formula for the first coefficient}
\label{sec:first}

\subsection{Linearity from the multivariate counting polynomial}
Equation~\eqref{eq:dilation} shows that $e_{n,1}(\bb)$ is the homogeneous
linear part of $F_n(\qq)$ evaluated at
$\qq=(b_1-1,b_2,\ldots,b_n)$. Hence it is enough to evaluate that linear
form on the $n$ coordinate axes. This observation avoids invoking a general
mixed-Ehrhart theorem or assuming a positivity property of mixed volumes.

For $1\le r\le n$, define the uniform-rank polytope
\begin{equation}
 Q_{r,n}=\{z\in[0,1]^n:z_1+\cdots+z_n\le r\}.
 \label{eq:uniform}
\end{equation}
On the $q_1$ axis, the rank polytope is a cube $q_1Q_{n,n}$. On the
$q_j$ axis for $j\ge2$, its $u$-sequence has $j-1$ zeros followed by
$n-j+1$ copies of $q_j$, so it is $q_jQ_{n-j+1,n}$.
These statements follow directly from the subset inequalities, including
their weak-parameter interpretation.

\subsection{The linear Ehrhart coefficient of a uniform-rank polytope}
\begin{lemma}
\label{lem:uniform-linear}
For $1\le r\le n$,
\begin{equation}
 [t]E_{Q_{r,n}}(t)=r(1+H_n-H_r).
 \label{eq:uniform-linear}
\end{equation}
\end{lemma}
\begin{proof}
At dilation $t$, count integer vectors $0\le z_i\le t$ with
$\sum_i z_i\le rt$. Inclusion--exclusion for the upper coordinate bounds
and the stars-and-bars formula give the polynomial identity
\begin{equation}
 E_{Q_{r,n}}(t)=\sum_{j=0}^{r-1}(-1)^j\binom nj
                     \binom{(r-j)t+n-j}{n}.
 \label{eq:uniform-ehrhart}
\end{equation}
For positive integer $t$, intersections of $j\ge r$ forbidden events
are empty, so they are not included. The terms shown with an insufficient
nonnegative upper binomial argument vanish in the usual combinatorial way.
Consequently \eqref{eq:uniform-ehrhart} is valid for every positive integer
$t$ and is a polynomial identity.

The derivative of the $j=0$ term at zero is $rH_n$, since
\[
 \binom{rt+n}{n}=\prod_{k=1}^n\left(1+\frac{rt}{k}\right).
\]
For $1\le j\le r-1$, the binomial polynomial has one zero factor at
$t=0$. Differentiating that factor gives
\[
 \left.\frac{d}{dt}\binom{(r-j)t+n-j}{n}\right|_{t=0}
 =\frac{(r-j)(n-j)!(-1)^{j-1}(j-1)!}{n!}.
\]
After multiplication by $(-1)^j\binom nj$, its contribution is
$-(r-j)/j$. Therefore
\[
 [t]E_{Q_{r,n}}=rH_n-\sum_{j=1}^{r-1}\frac{r-j}{j}
             =r(1+H_n-H_r).
\]
The formula includes $r=n$, where $Q_{n,n}$ is the unit cube and the
answer is $n$.
\end{proof}

\begin{proof}[Proof of Theorem~\ref{thm:first-intro}]
Coordinate-axis evaluation of the linear form gives
\begin{equation}
 e_{n,1}(\bb)=n(b_1-1)+
   \sum_{r=1}^{n-1}r(1+H_n-H_r)b_{n-r+1}.
 \label{eq:linear-ehrhart}
\end{equation}
Since $e_{n,0}=1$, the triangular transform gives
$\mu_{n,1}=e_{n,1}-n$. Substitute $b_i=1+x_i$ and evaluate the constant:
\begin{align*}
 \sum_{r=1}^{n-1}r(1+H_n-H_r)
 &=\frac{n(n-1)}2+
   \sum_{k=2}^n\frac1k\sum_{r=1}^{k-1}r\\
 &=\frac{n(n-1)}2+\frac{n(n-1)}4
 =\frac{3n(n-1)}4.
\end{align*}
This proves \eqref{eq:first-intro}. Every coefficient of an $x_i$ is
strictly positive. Thus the minimum on $x_i\ge0$ is attained uniquely at
$\xx=\zero$, where it equals $n(3n-7)/4$.
\end{proof}

\begin{remark}
At the extremes, the coefficients of $b_2$ and $b_n$ in
\eqref{eq:linear-ehrhart} are $(n^2-1)/n$ and $H_n$, respectively.
For $n=3$, these are $8/3$ and $11/6$, reproducing the asymmetric weights
in \eqref{eq:cubic-mu1}. At $n=2$, the constant minimum becomes $-1/2$,
exactly the exceptional magic coefficient above.
\end{remark}

\section{Endpoints, consequences, and parameter monotonicity}
\label{sec:consequences}

\subsection{The first and last entries of the magic vector}
\begin{proposition}
\label{prop:endpoints}
For all positive parameters and $n\ge3$,
\[
 \mu_{n,0}(\bb)=1,\qquad
 \mu_{n,n}(\bb)=\#\bigl(\interior\X_n(\bb)\cap\Z^n\bigr).
\]
Moreover $\mu_{n,n}(\bb)=0$ if and only if
$n=3$ and $\bb=(1,1,1)$.
\end{proposition}
\begin{proof}
The constant coefficient follows from $E(0)=1$. At $t=-1$, every term in
\eqref{eq:magic-basis} vanishes except the last, so
$\mu_{n,n}=(-1)^nE(-1)$. Ehrhart--Macdonald reciprocity identifies this with
the interior lattice count; this standard theorem is an external input
\cite{Macdonald,BeckDevelin}.

At $\bb=\one$, the original-coordinate upper bound for a subset of size
$k$ is
\[
 \Sigma_k=\frac{k(2n-k+1)}2.
\]
For $n\ge4$ this is strictly greater than $2k$ for every $1\le k\le n$.
Thus $(2,\ldots,2)$ strictly satisfies every lower and upper inequality
and is an interior lattice point. Increasing any parameter preserves
these strict inequalities.

For $n=3$ at $\bb=\one$, every interior lattice point would have all
coordinates at least two and total sum strictly less than six, which is
impossible. If any parameter is increased, the total-sum upper bound
becomes greater than six, while the bounds on one- and two-coordinate
sums were already strictly greater than two and four. Hence $(2,2,2)$
becomes interior. This proves the exact zero classification.
\end{proof}

\subsection{Consequences of the finite-dimensional theorem}
For $3\le n\le10$, the coefficientwise assertion has implications stronger
than nonnegativity at integer parameters. The polynomial continuations
$\mu_{n,j}(\one+\xx)$ are coordinatewise nondecreasing for real $x_i\ge0$,
and for $j\ge1$ they are strictly increasing in each coordinate. Their
minimum is therefore attained at the classical vector $\bb=\one$.
For $4\le n\le10$, every mixed partial derivative of total order at most $j$
is strictly positive on this orthant; the full-support statement proves
this by differentiating its matching monomial. These are statements about
the polynomial continuations, not an assertion that an irrational polytope
has an Ehrhart polynomial.

Since $\mu_{n,0}=1$, expanding the nonnegative magic basis gives
\begin{equation}
 [t^k]E_{n,\bb}(t)\ge\binom nk>0
 \qquad(0\le k\le n,\;3\le n\le10).
 \label{eq:ordinary-lower}
\end{equation}
The comparison is with $(1+t)^n$ coefficient by coefficient.

Define the $h^*$-polynomial by
\[
 \sum_{t\ge0}E_{n,\bb}(t)z^t=\frac{h^*_{n,\bb}(z)}{(1-z)^{n+1}}.
\]
Let $\mathcal E$ be the linear map sending $\binom tk$ to $w^k$.
The identity $\sum_{t\ge0}\binom tkz^t=z^k/(1-z)^{k+1}$ gives
\[
 h^*_{n,\bb}(z)=(1-z)^n
       \mathcal E(E_{n,\bb})\!\left(\frac{z}{1-z}\right).
\]
Br\"and\'en's theorem says that nonnegative magic coefficients force
$\mathcal E(E)$ to have all its zeros in $[-1,0]$
\cite[Theorem~4.2 and equation~(6.4)]{Branden}. The fractional-linear
substitution takes these zeros to nonpositive real zeros; a zero at $-1$
accounts for a possible degree drop. Thus the finite-dimensional theorem
also gives real-rooted $h^*$-polynomials for every positive parameter vector
with $3\le n\le10$. Their coefficient sequences are consequently
log-concave and unimodal. The real-rootedness implication is classical;
only its application to the newly certified parameter range is part of
the present result.

\section{Cube addition and a reduction to the boundary parameter}
\label{sec:cube}

The following identity clarifies which parameter is easiest to remove from
a possible all-dimensional proof. It is an exact identity, not a positivity
assumption about general Minkowski sums.

\begin{lemma}[Adding coordinate segments]
\label{lem:segments}
Let $P\subset\R_{\ge0}^n$ be a compact downward-closed convex set. For an
integer $m\ge0$ and a coordinate vector $e_i$,
\[
 \#((P+[0,m]e_i)\cap\Z^n)
 =\#(P\cap\Z^n)+m\#(\proj_{[n]\setminus\{i\}}P\cap\Z^{n-1}).
\]
\end{lemma}
\begin{proof}
Over an integer point of the projection, the original fiber is a nonempty
interval $[0,u]$. It contains $\lfloor u\rfloor+1$ integers. After adding
the segment it is $[0,u+m]$, containing exactly $m$ more. No other fibers
appear. Downward closure is what guarantees that a nonempty projected
fiber contains the integer height zero.
\end{proof}

By the greedy support description,
\begin{equation}
 P_n(b_1+a,b_2,\ldots,b_n)=P_n(\bb)+a[0,1]^n
 \qquad(a\ge0).
 \label{eq:cube-sum}
\end{equation}
For $1\le k\le n$, a projection onto any $k$ coordinates has the same
rank bounds through size $k$. Its parameter vector is
\[
 \pi_k\bb=(S_{n-k+1},b_{n-k+2},\ldots,b_n).
\]
Put $E_{0,\pi_0\bb}=M_{0,\pi_0\bb}=1$ as an empty-dimensional convention.

\begin{proposition}[Cube-addition formula]
\label{prop:cube}
For nonnegative integer $a$ and positive $\bb$,
\begin{align}
 E_{n,(b_1+a,b_2,\ldots,b_n)}(t)
 &=\sum_{k=0}^n\binom nk (at)^{n-k}E_{k,\pi_k\bb}(t),
 \label{eq:cube-ehrhart}\\
 M_{n,(b_1+a,b_2,\ldots,b_n)}(y)
 &=\sum_{k=0}^n\binom nk (ay)^{n-k}M_{k,\pi_k\bb}(y).
 \label{eq:cube-magic}
\end{align}
The degree-$k$ ambient transform is used on the right; the projected
one-dimensional parameters occurring for $k<n$ are at least two.
\end{proposition}
\begin{proof}
At an integer dilation $t$, apply Lemma~\ref{lem:segments} successively
in all coordinates, with segment length $at$. Choosing the coordinates
in which the projection term is taken produces the sum over all subsets.
Symmetry identifies all terms with the same remaining dimension $k$,
giving \eqref{eq:cube-ehrhart}. Projection commutes with addition of the
other coordinate segments, so no mixed term is omitted.

Both sides are polynomials in $t$. Substituting $t=y/(1-y)$ and multiplying
by $(1-y)^n$ turns each term into
$\binom nk(ay)^{n-k}M_{k,\pi_k\bb}(y)$, which proves
\eqref{eq:cube-magic}.
\end{proof}

\begin{corollary}[Boundary reduction]
\label{cor:boundary}
To prove Conjecture~\ref{conj:target} in every dimension, it is sufficient
to prove it for the parameter vectors $(1,b_2,\ldots,b_n)$ in every
dimension $n\ge3$.
\end{corollary}
\begin{proof}
Use induction on $n$, starting with the elementary classifications in
dimensions one and two. For a desired vector, apply
\eqref{eq:cube-magic} with base vector $(1,b_2,\ldots,b_n)$ and
$a=b_1-1$. The term $k=n$ is nonnegative by the boundary hypothesis.
Terms $3\le k<n$ are nonnegative by induction. For $k=2<n$, the first
projected parameter is at least $n-1\ge2$, excluding the only
non-magic-positive two-dimensional vector $(1,1)$. For $k=1<n$, it is
at least $n\ge3$, so its magic transform is nonnegative. The term $k=0$
is nonnegative as well.
\end{proof}

This is a reduction of the difficulty, not its disappearance: the boundary
family still has $n-1$ independent positive parameters. Nor does the formula
say that arbitrary Minkowski addition preserves magic positivity. Its
special form depends on coordinate segments and downward-closed fibers.

\section{Verification, reproducibility, and the limits of the attempt}
\label{sec:verification}

\subsection{Independent checks performed}
The recorded run of \texttt{code/verify\_certificates.py} completed with
status \texttt{PASS}. It performs full coefficient-by-coefficient
regeneration of all lattice and magic certificates through dimension ten,
checks the degree and support claims, and compares every first-coefficient
map with the independent harmonic-number formula. It also verifies all 55
kernel difference identities with $a+b\le9$ by direct polynomial
translation.

The independent geometric enumerator uses only
\eqref{eq:rank-polytope}. It lists decreasing nonnegative integer tuples
$z_1\ge\cdots\ge z_n$ and tests
\[
 z_1+\cdots+z_k\le t f(k)\qquad(1\le k\le n).
\]
These are equivalent to every subset inequality, since the largest sum
of $k$ coordinates is the prefix sum. A tuple with multiplicities
$m_1,m_2,\ldots$ represents exactly $n!/\prod_i m_i!$ coordinate
permutations. Summing these orbit sizes gives the full lattice count.
This procedure uses neither the slice recurrence nor power sums nor
Ehrhart interpolation.

\begin{center}
\begin{tabular}{@{}llllr@{}}
\toprule
$n$ & Parameter vectors & Dilations & & Counts\\
\midrule
3 & $\{1,2,3,4\}^3$ & $0,1,2,3$ &&256\\
4 & $\{1,2,3\}^4$ & $0,1,2$ &&243\\
5 & $\{1,2\}^5$ & $0,1,2$ &&96\\
6 & $\{1,2\}^6$ & $0,1$ &&128\\
7--10 & $\bb=\one$ & $0,1$ &&8\\
\midrule
\multicolumn{4}{l}{Total} &731\\
\bottomrule
\end{tabular}
\end{center}
There are 486 positive-dilation checks and 245 zero-dilation checks. Each
agrees both with the lattice polynomial under \eqref{eq:dilation} and
with the magic-basis expansion. All inputs and resulting integer counts
are recorded in \texttt{data/verification.json}. These finite geometric
checks are error-detection tests; the all-parameter conclusion comes
from the symbolic identities and nonnegative coefficients, not from
extrapolating the tests.

\subsection{Reproducing and using the package}
The core programs use only the Python 3.9+ standard library. From the
unpacked package directory:
\begin{lstlisting}
python code/verify_certificates.py
python code/evaluate_polytope.py 2 3 1
python code/sparse_magic.py --max-dimension 10 --output data/rebuilt
pdflatex -interaction=nonstopmode -halt-on-error article.tex
pdflatex -interaction=nonstopmode -halt-on-error article.tex
\end{lstlisting}
The first command regenerates the full certificates before checking them.
The evaluator returns exact rational Ehrhart and magic coefficients,
integer $h^*$-coefficients, and a requested lattice count. It reads the
provided certificates and accepts dimensions two through ten. Higher
symbolic dimensions can be attempted with the generator, but are not
part of the theorem certified by this archive.

The supplied run took approximately 33 seconds for full verification in
the execution environment used here; this is a recorded observation,
not a hardware-independent performance promise. The uncompressed
certificate directory is about 25 MB. Its largest files are sparse
coefficient lists, not opaque binary proof objects.

\subsection{Why the natural induction has not yet closed}
The data suggest a stronger all-dimensional assertion:
\begin{conjecture}[Coefficientwise strengthening]
\label{conj:strong}
For every $n\ge3$ and $0\le j\le n$,
$\mu_{n,j}(\one+\xx)$ has nonnegative ordinary coefficients.
For $n\ge4$ its support is the entire set of monomials of degree at most $j$.
\end{conjecture}
The numerical dimension bound has not been removed. In particular, the
positive coefficients seen through dimension ten do not justify an
induction unless the recursion is shown to preserve the relevant cone
of polynomials.

A specific obstacle is that the power-sum operation does not preserve
ordinary coefficientwise nonnegativity. Already
\[
 \sum_{r=1}^Y r^4
 =\frac{Y^5}{5}+\frac{Y^4}{2}+\frac{Y^3}{3}-\frac{Y}{30}
\]
has a negative coefficient. The magic transform also contains alternating
binomial signs. The actual lattice polynomials exhibit cancellations
that make the final shifted magic coefficients positive in the certified
range, but an arbitrary nonnegative input polynomial need not do so.
A proof that merely calls each summation step ``positive'' would be invalid.

The attempts here isolate three concrete routes for further work.
First, it would suffice to build a positivity-preserving recurrence on the
boundary family $b_1=1$, using Corollary~\ref{cor:boundary} to restore the
first parameter. Second, a combinatorial interpretation for the integers
$[\xx^\alpha]n!\mu_{n,j}(\one+\xx)$ would explain the full support and
remove the dimension cutoff simultaneously. Third, one can seek an
invariant stronger than ordinary coefficientwise positivity that is
preserved by the finite-sum kernels together with the triangular transform.
The published two-parameter theorem remains a useful exact specialization
against which any candidate recurrence must be checked.

\subsection{Status of the mathematical claims}
The theorem in dimensions three through ten is presented as a
computer-assisted mathematical result with explicit proof dependencies and
reproducible exact data. The cubic formulas, the all-dimensional first
coefficient, the endpoint classification, and the cube identity have
ordinary proofs printed above. The inequality description and slice
machinery are credited to their sources; reciprocity and Br\"and\'en's
theorem are stated external inputs. Neither global magic positivity nor
global coefficientwise positivity is claimed. The draft has not been
independently refereed or checked in a proof assistant, and priority for
the consequences derived here has not been established independently.

\appendix
\section{Random selection and the full area catalogue}
\label{app:selection}
The generator \texttt{code/select\_area.py} first fixes the ordered list
below, obtains a seed using \texttt{secrets.randbits(128)}, and calls
\texttt{random.Random(seed).randrange(120)}. The recorded seed is
\[
 \texttt{89054884759964478034823739394527898187}.
\]
It returns zero-based index 32, or one-based entry 33. To replay the same
draw without obtaining new entropy:
\begin{lstlisting}
python code/select_area.py --seed 89054884759964478034823739394527898187
\end{lstlisting}
The selection is reproducible, but its purpose is diversification rather
than a mathematical claim about a canonical probability distribution on
research areas. In particular, the list itself is an editorial choice.

\begin{multicols}{2}
\small
\begin{enumerate}[label=\arabic*.,itemsep=0pt,topsep=0pt,leftmargin=2.7em]
\item Additive combinatorics
\item Algebraic combinatorics
\item Algebraic geometry
\item Algebraic graph theory
\item Algebraic number theory
\item Algebraic topology
\item Analytic combinatorics
\item Analytic number theory
\item Approximation theory
\item Automata theory
\item Banach space theory
\item Braid groups
\item Calculus of variations
\item Cellular automata
\item Coding theory
\item Combinatorial commutative algebra
\item Combinatorial designs
\item Combinatorial game theory
\item Combinatorial geometry
\item Combinatorics on words
\item Commutative algebra
\item Complex dynamics
\item Computability theory
\item Computational geometry
\item Convex geometry
\item Cryptographic Boolean functions
\item Descriptive set theory
\item Diophantine approximation
\item Difference equations
\item Differential algebra
\item Differential geometry
\item Discrete dynamical systems
\item \textbf{Discrete geometry}
\item Discrepancy theory
\item Distance geometry
\item Enumerative combinatorics
\item Ergodic theory
\item Extremal graph theory
\item Finite fields
\item Finite group theory
\item Formal languages and grammars
\item Fourier analysis
\item Functional equations
\item General topology
\item Geometric group theory
\item Geometry of numbers
\item Graph colorings
\item Graph polynomials
\item Harmonic analysis
\item Hypergraph theory
\item Incidence geometry
\item Information theory
\item Integer partitions
\item Integral transforms
\item Invariant theory
\item Knot theory
\item Lattice theory
\item Lie algebras
\item Linear algebra and matrix inequalities
\item Logic and proof theory
\item Low-dimensional topology
\item Mathematical billiards
\item Mathematical probability
\item Matroid theory
\item Metric geometry
\item Model theory
\item Modular forms
\item Moment problems
\item Nonassociative algebra
\item Noncommutative algebra
\item Numerical analysis
\item Operator theory
\item Optimization
\item Ordered algebraic structures
\item Order theory
\item Orthogonal polynomials
\item Packing and covering
\item Partial differential equations
\item Permutation patterns
\item Polynomial inequalities
\item Polyhedral combinatorics
\item Potential theory
\item q-Series and basic hypergeometric functions
\item Quiver representations
\item Ramsey theory
\item Random graphs
\item Real algebraic geometry
\item Recurrence sequences
\item Representation theory
\item Riordan arrays
\item Semigroup theory
\item Set theory
\item Spectral graph theory
\item Special functions
\item Substitution dynamical systems
\item Symbolic dynamics
\item Symmetric functions
\item Tensor decomposition
\item Topological dynamics
\item Tropical geometry
\item Type theory
\item Ultrametric analysis
\item Uniform distribution
\item Universal algebra
\item Variational inequalities
\item Vertex algebras
\item Well-quasi-orders
\item Zero-sum theory
\item Discrete differential geometry
\item Algebraic statistics
\item Numerical semigroups
\item Positivity and total positivity
\item Rewriting systems
\item Computational number theory
\item Fractional calculus
\item Nonstandard analysis
\item Geometric measure theory
\item Boolean lattices and set systems
\item Random walks
\item Higher category theory
\end{enumerate}

\end{multicols}

\section{Certificate format and audit}
\label{app:format}
A lattice certificate named \texttt{lattice\_nn.json} stores
\texttt{dimension}, \texttt{scale}, and a lexicographically sorted
\texttt{terms} list. Each term is \texttt{[[a1,...,an],c]} and denotes
$c\prod_i b_i^{a_i}$ in the scaled polynomial $n!L_n$.
A file \texttt{magic\_nn\_jj.json} uses exactly the same representation
but denotes $n!\mu_{n,j}(\one+\xx)$ in the $x_i$ variables.
The distinction between the two variable conventions is essential.

For example, the dimension-three first magic certificate is the list of
terms
\[
 ((0,0,0),9),\quad ((0,0,1),11),\quad
 ((0,1,0),16),\quad ((1,0,0),18),
\]
with scale $6$. Dividing by six reproduces
\eqref{eq:cubic-mu1}. In the dimension-three final certificate, the
constant term is absent; its value is exactly zero, not a discarded small
floating-point value.

A minimal coefficient audit, independent of any symbolic package, is:
\begin{lstlisting}[language=Python]
import json
from math import comb, factorial
from pathlib import Path

root = Path("data/certificates")
for n in range(3, 11):
    for j in range(n + 1):
        d = json.loads((root / f"magic_{n:02d}_{j:02d}.json").read_text())
        assert d["dimension"] == n and d["scale"] == factorial(n)
        terms = d["terms"]
        exponents = [tuple(a) for a, c in terms]
        assert len(set(exponents)) == len(exponents)
        assert all(len(a) == n and min(a) >= 0 for a in exponents)
        assert all(sum(a) <= j for a in exponents)
        assert all(isinstance(c, int) and c > 0 for a, c in terms)
        assert len(terms) == comb(n + j, j) - ((n, j) == (3, 3))
\end{lstlisting}
This small audit establishes positivity of the supplied polynomial lists,
but by itself does not establish that they are the right polynomials.
For that, use the full verifier, which checks their exact construction
from the geometric recurrence. Keeping these two obligations separate
is part of the proof, not merely a software convention.

\begin{thebibliography}{9}
\bibitem{HLTV}
C.~Hill, A.~Luo, V.~Trinh, and A.~R.~Vindas-Mel\'endez,
\emph{Lattice slices, Ehrhart polynomials, and magic positivity of generalized
parking-function polytopes}, preprint, 2026, arXiv:2607.15503v1.
\url{https://arxiv.org/abs/2607.15503v1}.

\bibitem{Bayer}
M.~M.~Bayer, S.~Borgwardt, T.~Chambers, S.~Daugherty, A.~Dawkins,
D.~Deligeorgaki, H.-C.~Liao, T.~McAllister, A.~Morrison, G.~Nelson,
and A.~R.~Vindas-Mel\'endez,
\emph{Combinatorics of generalized parking-function polytopes},
Discrete Comput. Geom. \textbf{76} (2026), 339--377.
\url{https://doi.org/10.1007/s00454-025-00770-1}.
Preprint: \url{https://arxiv.org/abs/2403.07387}.

\bibitem{Branden}
P.~Br\"and\'en, \emph{On linear transformations preserving the P\'olya
frequency property}, Trans. Amer. Math. Soc. \textbf{358} (2006),
3697--3716. Preprint: \url{https://arxiv.org/abs/math/0403364}.

\bibitem{Macdonald}
I.~G.~Macdonald, \emph{Polynomials associated with finite cell-complexes},
J. London Math. Soc. (2) \textbf{4} (1971), 181--192.
\url{https://doi.org/10.1112/jlms/s2-4.1.181}.

\bibitem{BeckDevelin}
M.~Beck and M.~Develin, \emph{On Stanley's reciprocity theorem for rational
cones}, preprint, 2004, arXiv:math/0409562.
\url{https://arxiv.org/abs/math/0409562}.
This supplies an accessible proof of a reciprocity theorem implying the
Ehrhart--Macdonald identity used here.
\end{thebibliography}
\end{document}
